\documentclass[10pt,a4paper]{amsart}
\usepackage[margin=19mm]{geometry}
\usepackage{amsmath,amssymb,mathtools}
\usepackage[scr=rsfs,cal=euler]{mathalfa}
\usepackage{xfrac}
\usepackage[unicode,hidelinks,bookmarks=false]{hyperref}
\newcommand{\ie}{i.e.\ }
\newcommand{\cf}{cf.\ }
\newcommand{\resp}{respectively\ }
\newcommand{\Subsection}{\subsection}
\providecommand{\nobreakdash}{\nobreak\hbox{-}}
\setlength{\emergencystretch}{2em}
\multlinegap0pt
\usepackage[shortlabels]{enumitem}
\usepackage{amssymb}
\usepackage{mathtools}
\newtheorem{theorem}{Theorem}
\newcommand{\F}{\mathbb{F}}
\newcommand{\Ocal}{\mathcal{O}}
\newcommand{\Q}{\mathbb{Q}}
\newcommand{\Z}{\mathbb{Z}}
\newcommand{\disc}{\operatorname{disc}}
\title{Eisenstein-prime Obstruction Sieve for Monogenicity}
\author{Khai-Hoan Nguyen-Dang}
\date{Source: arXiv:2602.10488v1 (CC BY 4.0)}
\begin{document}
\maketitle

\noindent\emph{Source and licence.} Eisenstein-prime Obstruction Sieve for Monogenicity. Version arXiv:2602.10488v1 (CC BY 4.0), distributed by arXiv under the Creative Commons Attribution 4.0 International licence, http://creativecommons.org/licenses/by/4.0/, as recorded in the arXiv licence metadata for this exact version and shown on the abstract page https://arxiv.org/abs/2602.10488v1. Full licence notice: \texttt{licenses/arxiv-2602.10488-CC-BY-4.0.txt}.\par\smallskip

\noindent\textit{Editorial scope.} Counted unit: the article's theorem thm:ft-monogenic-positive-density (source0.tex lines 2080--2088). The article's complete original proof of it is delivered with it (source0.tex lines 2090--2164, one \texttt{proof} environment of 75 lines). No mathematics is written, repaired or re-derived: the statement and the proof are the article's own lines, in the article's own notation, and no retained line is substituted or re-flowed.  Retained uncounted as the source-local context the proof needs: lines 2064--2067.  Nothing was omitted from any retained span, so the package carries no omission record.  No retained line contains a citation, so the wrapper carries no bibliography and no bibliography entry.  No retained line contains a figure environment, an image-inclusion command or drawing code, so no image is delivered and the package declares no asset dependency.  Editorial formatting only, in addition: an AMS \texttt{amsart} preamble that restates the article's own declarations and macros used by the retained lines, namely the environments \texttt{theorem} on separate counters, together with the standard packages those lines need.  The retained environments print in this document as Theorem 1 in order of appearance, whereas the article prints the same items with the numbering of its own full document.  The complete native source of the article as distributed in the arXiv e-print for this exact version is bundled unchanged as \texttt{sources/194402/source0.tex}.
\begin{equation}
\label{eq:disc-ft}
\disc(f_t)=t^{n-1}L_n(t).
\end{equation}

\begin{theorem}[Positive density of monogenic (hence \textsc{ABS}--unobstructed) specializations]
\label{thm:ft-monogenic-positive-density}
For every $n\ge 4$, the set $\mathcal T_n$ has positive natural density in $\Z$.
Moreover, for every $t\in\mathcal T_n$, the order $\Z[\alpha_t]$ is maximal:
\[
\Ocal_{K_t}=\Z[\alpha_t].
\]
In particular, $K_t$ is monogenic and hence \textsc{ABS}--unobstructed for all $t\in\mathcal T_n$.
\end{theorem}

\begin{proof}
Fix $t\in\mathcal T_n$ and put $I_t:=[\Ocal_{K_t}:\Z[\alpha_t]]$.

Let $q\mid t$ be prime. Since $t$ is squarefree and $\gcd(t,n(n-1))=1$, we have $v_q(t)=1$ and $q\nmid n(n-1)$.
The polynomial $f_t(X)=X^n+tX+t$ is Eisenstein at $q$, hence irreducible over $\Q$.


Fix a prime $q\mid t$ and set
\[
R_q:=\Z_q[\alpha_t]\;\cong\;\Z_q[X]/(f_t).
\]
Since $f_t$ is Eisenstein at $q$, it is irreducible over $\Q_q$, hence $R_q$ is a domain, finite over the DVR $\Z_q$,
and therefore $\dim(R_q)=1$. Reducing modulo $q$ gives
\[
R_q/qR_q \cong \F_q[X]/(X^n),
\]
so $R_q$ is local with maximal ideal $\mathfrak m_q=(q,\alpha_t)$.

In $R_q$ we have $\alpha_t^n=-t(\alpha_t+1)$, and since $\alpha_t\in\mathfrak m_q$ we have
$\alpha_t+1\equiv 1\pmod{\mathfrak m_q}$, hence $\alpha_t+1\in R_q^\times$.
Thus $t\in(\alpha_t)$, and writing $t=qu$ with $u\in\Z_q^\times$ gives $q\in(\alpha_t)$.
Consequently $\mathfrak m_q=(q,\alpha_t)=(\alpha_t)$ is principal.

Since $R_q$ is a $1$-dimensional Noetherian local domain with principal maximal ideal, it is a DVR.
In particular $R_q$ is integrally closed, hence equals the ring of integers of $\Q_q(\alpha_t)$.
Therefore $\Ocal_{K_t}\otimes\Z_q \cong R_q$, so $q\nmid I_t$.


If $p\mid I_t$ then $p^2\mid \disc(f_t)$ by the index--discriminant relation
$\disc(f_t)=\disc(K_t)\,I_t^2$.
If $p\nmid t$, then \eqref{eq:disc-ft} implies $p^2\mid L_n(t)$.
But $tL_n(t)$ is squarefree by hypothesis, hence $L_n(t)$ is squarefree; contradiction.
Thus $p\nmid I_t$ for all primes $p\nmid t$.

Combining above observations yields $I_t=1$, hence $\Ocal_{K_t}=\Z[\alpha_t]$ and $K_t$ is monogenic.


Write
\[
F(t):=t\,L_n(t)=t\,(C_{0,n}+C_{1,n}t)\in\Z[t].
\]
Since $\gcd(C_{0,n},C_{1,n})=1$, the two linear factors $t$ and $L_n(t)$ are coprime in $\Z[t]$,
so $F$ is a squarefree polynomial.

For each prime $\ell$, let
\[
\rho(\ell^2):=\#\{\,a\in\Z/\ell^2\Z:\ \ell^2\mid F(a)\,\}.
\]
A direct congruence count gives:
\begin{enumerate}[label=\textup{(\roman*)}]
\item If $\ell\mid n$ (equivalently $\ell\mid C_{0,n}$), then $C_{1,n}\not\equiv 0\pmod\ell$ and
$L_n(a)\equiv C_{1,n}a\pmod\ell$. Hence $\ell^2\mid aL_n(a)$ iff $\ell\mid a$, so $\rho(\ell^2)=\ell$.
\item If $\ell\mid (n-1)$ (equivalently $\ell\mid C_{1,n}$), then $C_{0,n}\not\equiv 0\pmod\ell$, so
$L_n(a)\not\equiv 0\pmod\ell$ for all $a$, and therefore $\ell^2\mid aL_n(a)$ iff $\ell^2\mid a$.
Thus $\rho(\ell^2)=1$.
\item If $\ell\nmid n(n-1)$, then $C_{0,n}$ and $C_{1,n}$ are both units modulo $\ell^2$ and the two linear
congruences $a\equiv 0\pmod{\ell^2}$ and $L_n(a)\equiv 0\pmod{\ell^2}$ each contribute exactly one solution.
Since they are distinct (because $C_{0,n}\not\equiv 0\pmod\ell$), we have $\rho(\ell^2)=2$.
\end{enumerate}

In particular, for every prime $\ell$ we have $\rho(\ell^2)<\ell^2$, and the Euler product
\[
\mathfrak S:=\prod_{\ell}\Bigl(1-\frac{\rho(\ell^2)}{\ell^2}\Bigr)
\]
converges absolutely to a positive real number (since $\rho(\ell^2)\ll 1$ for all but finitely many $\ell$,
and $\sum_\ell \ell^{-2}<\infty$). By the standard squarefree sieve/inclusion--exclusion for polynomial values
in one variable (applied to $F(t)$), the set
\[
\{t\in\Z:\ F(t)\ \text{is squarefree}\}
\]
has natural density $\mathfrak S>0$. Intersecting with the congruence condition $\gcd(t,n(n-1))=1$
removes only finitely many residue classes modulo $\prod_{\ell\mid n(n-1)}\ell$, hence preserves positivity
of density. Finally, excluding $|t|\le 1$ removes finitely many integers. Therefore $\mathrm{dens}(\mathcal T_n)>0$.

\end{proof}

\end{document}
